\section{引言}
\label{sec:introduction}

在 UTXO 账本中，每笔支付都花费先前交易的输出，因此若接收方等待公开执行，支付链的每一跳都会增加一次确认延迟。假设 Alice 向 Bob 支付 100 CAL，而 Bob 在 Alice 的支付执行之前将其转给 Carol。若 Bob 的支付先执行，账本将消费一个尚不存在的输出，缺失的 100 CAL 就需要一个明确且有资金支持的来源，无论 Alice 的交易稍后到达还是永不到达。

在每一跳都等待公开执行是最直接的替代方案，我们的评估将其用作对照。基于证书的系统把接收方的可用性与公开处理分离开来：FastPay 和 Zef 支持低延迟的认证支付，Sui Lutris 则针对不同类型的状态访问结合了广播路径与共识路径~\cite{fastpay,zef,lutris}。支付通道依靠通道流动性实现快速转账~\cite{blitz,donner}。Next 要问的是：原生账本如何为来源尚未执行的认证支付提供资金，其背后由已注册的组织而非通道作为支撑。

三个困难相互耦合。证书在公开登记之前就已流通，因此仅依据公开缺口进行核算会遗漏已存在的承诺。迟到的来源交易不得既向垫付价值的担保方付款，又向已经花掉该价值的接收方付款。客户端可能持有某个来源交易却永不提交。

我们的核心观察是：对缺失来源负责的一方，已经由随输出流转的证书指明。Next 将缺口记在该直接发行方名下，而不是追溯祖先交易；以由备付金支持的签名额度约束每份已签证书的责任；并通过公开决定关闭未解决的义务，而不改变已存储的历史。

经济闭合之后，还有一个独立的记录问题：已赔付输入在原始区块中仍引用缺失输出。账本审计已有利用保存的执行证据回访历史记录的工作流~\cite{ia-ccf}；Reparo 也通过原交易哈希返回修复后的交易数据~\cite{reparo}。我们将这一工作流用于赔付：归档副本可在原付款坐标处导出修订后的资金引用。索引能够提供相同的经济答案。新增的技术任务是授权替换正文，同时保持既有提交证据和后继区块头所引用的完整身份~\cite{comet-blockid}，且不重新执行该支付。

该设计包含三项相互补充的贡献。

\noindent\textbf{C1：支持一轮继续的分层架构。}认证与公开执行承担不同的职责：证书固定输出并将责任归于其发行方，而公开执行决定账本接受哪笔支付。消费组织的成员在签名前原子地锁定输入并预留额度，一轮并行交互即产生一个 TXCer。接收方验证并记录该输出后可立即请求后继交易，且只携带其直接输入的证书，复制与公开提交则在后台运行。后继批准仍须通过输入、费用和容量检查；第~\ref{sec:model-contract}~节给出分阶段的服务契约。

\noindent\textbf{C2：带来源解耦的直接责任。}公开执行以原子方式消费缺失的已认证输入，并针对其发行方记录一项义务。后继交易所在区块公布来源交易的证书，因此原签名者可以根据存储的批准重新提交该来源交易；截止期之后，公开决定无需门限适配即可扣减发行方的备付金，迟到的来源交易则将其偿还。每个批准在仍可能被认证期间预留完整的 CAL 输出（含找零），而经认证的冲突性公开花费可以释放未获证批准的未使用本地预留。对于由四名成员组成的组织和由四个验证者组成的委员会（各自至多含一个静态拜占庭成员），我们证明了唯一消费，以及签名额度覆盖每份证书（无论是否已公布）未偿 CAL 责任这一性质。对于一个有限、无环、无冲突、包含其已认证输入之来源交易且每笔支付只执行一次的支付集合，来源到达与赔付的所有顺序最终得到相同的 CAL 持有量和备付余额。该核算将成员间重叠的预算~\cite{stingray}扩展到隐藏的证书责任、动态赔付与偿还以及来源闭合。

\noindent\textbf{C3：由赔付决定授权的资金引用修复。}已提交的赔付决定在经济上关闭义务，C3 则将该决定落实到存储的历史中。受限变色龙哈希命令将其备付扣款引用写入指定输入。精确的修订与任务授权这一替换，同时保持所有者授权、输出及后继引用和完整 BlockID 不变。原始命令保留用于重放，表示操作不写入核算状态。这支持在既有区块身份下进行本地异常复核与修订区块再导出，而不让经济闭合等待历史修复。

在一台 M4 Max 主机上，采用成员同步提交和 NoSync 钱包存储时，
100 跳链到达最后一笔认证到账的时长中位数为 4.707\,s，
跳间等待公共执行时则为 64.479\,s。
同一当前构建的三轮十四服务 P 实验共完成 21,048 笔付款；
赔付和回款均在暂停的适配服务恢复前完成。
独立实验还评估注入网络延迟和单一故障下的运行行为，
以及可选历史表示的本地读取、存储和维护成本。
